% =============================================================================
% CHAPTER 8: Computational Tools, Roadmap & References
% Sources: tools.md, part4_development_roadmap.md, about.md
% =============================================================================

\chapter{Computational Tools, Roadmap \& References}
\label{ch:tools}

% -----------------------------------------------------------------------------
\section{Available Computational Scripts}
\label{sec:scripts}

We provide a suite of Python tools for exploring the oscillating brane theory and computing its predictions. All scripts are available at: \url{https://github.com/Teleadmin-ai/oscillating-brane-DM/tree/main/scripts}

\begin{table}[h]
\centering
\small
\caption{V8.0 computational scripts and their outputs.}
\label{tab:scripts}
\begin{tabular}{@{}p{4.2cm}p{6.5cm}p{3.5cm}@{}}
\toprule
\textbf{Script} & \textbf{Purpose} & \textbf{Output} \\
\midrule
\texttt{brane\_dynamics.py} & Core V8.0 ODE ($\Fweb + \Rpbh$ + trace + $\Grad$ + $\xi R\phi$) & Console \\
\texttt{radion\_attractor.py} & Dynamical attractor demonstration & \texttt{plots/radion\_attractor\_*.png} \\
\texttt{pbh\_emf\_constraints.py} & PBH EMF vs microlensing (wave-optics) & \texttt{plots/pbh\_emf\_constraints.png} \\
\texttt{bbn\_thermal\_freezeout.py} & BBN via conformal symmetry \& trace anomaly & \texttt{plots/bbn\_thermal\_freezeout.png} \\
\texttt{growth\_scale\_dependent.py} & Scale-dependent $S_8$ Yukawa screening & \texttt{plots/growth\_scale\_dependent.png} \\
\texttt{numerical\_relativity\_1d.py} & 5D radiative damping (1+1)D MoL & \texttt{plots/warped\_shielding\_1D.png} \\
\texttt{desi\_w\_evolution.py} & Dark energy $w(z)$ matching DESI & \texttt{plots/desi\_w\_evolution.png} \\
\texttt{s8\_tension\_resolution.py} & Structure growth suppression & \texttt{plots/s8\_tension\_resolution.png} \\
\texttt{isw\_cmb\_signature.py} & ISW resonance analysis & \texttt{plots/isw\_cmb\_signature.png} \\
\texttt{bayesian\_analysis.py} & MCMC parameter inference & Posterior chains \\
\texttt{kk\_spectrum.py} & Kaluza--Klein graviton spectrum & Console \\
\texttt{growth\_factor.py} & Growth factor $D_+(z)$ computations & Console \\
\bottomrule
\end{tabular}
\end{table}

% -----------------------------------------------------------------------------
\section{API Documentation}
\label{sec:api}

\paragraph{BraneOscillator Class.}
\begin{verbatim}
class BraneOscillator:
    def __init__(self, tau_0=7.0e19, f_osc=0.10, T=2.0, L=2.0e-7):
        """
        Parameters:
        - tau_0: Brane tension (J/m^2)
        - f_osc: Oscillating DM fraction
        - T:     Period (Gyr)
        - L:     Extra dimension size (m)
        """
\end{verbatim}

Key functions: \texttt{equation\_of\_state(z)}, \texttt{membrane\_displacement(t)}, \texttt{growth\_suppression()}.

\paragraph{GrowthFactorCalculator Class.}
\begin{verbatim}
class GrowthFactorCalculator:
    def __init__(self, omega_m=0.315, oscillating=True, A_w=0.003):
\end{verbatim}

Supports fast fitting formula or exact ODE integration (\texttt{-{}-exact} flag), comparison between oscillating and \LCDM{} models, and $S_8$ parameter calculation.

% -----------------------------------------------------------------------------
\section{Observational Tests Timeline 2025--2035+}
\label{sec:obs_timeline}

\paragraph{2025--2027 (Near Term):}
\begin{itemize}
  \item \textbf{Euclid}: Wide-field weak lensing $\to$ $S_8$ precision to 1\%.
  \item \textbf{DESI}: BAO measurements $\to$ $w(z)$ amplitude constraints.
  \item \textbf{CMB-S4}: Large-scale anisotropies $\to$ ISW resonance detection.
  \item \textbf{JWST}: Ultra-faint dwarf census $\to$ subhalo abundance.
\end{itemize}

\paragraph{2028--2030 (Medium Term):}
\begin{itemize}
  \item \textbf{Vera Rubin Observatory (LSST)}: 10-year survey $\to$ halo profiles to 200~kpc; stellar streams $\to$ substructure constraints; microlensing $\to$ smooth vs clumpy halos.
  \item \textbf{Roman Space Telescope}: High-$z$ structure $\to$ growth history.
  \item \textbf{SKA-Low}: 21\,cm reionization modulation $\to$ \textbf{definitive test}.
\end{itemize}

\paragraph{2030--2035 (Long Term):}
\begin{itemize}
  \item \textbf{CMB-S4 Full Survey}: Precision ISW measurements; $2\Gyr$ oscillation signature.
  \item \textbf{ELT/TMT}: Dwarf galaxy kinematics $\to$ core sizes.
  \item \textbf{Advanced gravitational tests}: $\delta g/g$ measurements.
\end{itemize}

\paragraph{2035+ (Future):}
\begin{itemize}
  \item Next-gen CMB experiments: Enhanced ISW detection.
  \item Next-gen atom interferometry: Spatial gravity variations.
  \item Combined CMB + galaxy surveys: Definitive oscillation mapping.
\end{itemize}

% -----------------------------------------------------------------------------
\section{24-Month Development Roadmap}
\label{sec:roadmap}

\paragraph{Phase 1: Advanced Theoretical Framework (Months 1--6).}
Formulate complete 5D action with Goldberger--Wise stabilization. Implement ADM (3+1)+1 decomposition for numerical evolution. Derive Israel junction conditions for oscillating brane boundary. Choose optimal gauge (Gaussian normal or Eddington--Finkelstein coordinates).

\paragraph{Phase 2: Numerical Infrastructure (Months 6--12).}
Extend GRChombo to 5D geometry with adaptive mesh refinement. Implement moving boundary conditions with $\mathbb{Z}_2$ symmetry. Develop Python/Julia prototypes for rapid testing. Validate against known static solutions (RS metric recovery).

\paragraph{Phase 3: Physical Applications (Months 12--18).}
Simulate brane collision scenarios ($v_\text{rel} \sim 10^{-3}c$). Include inflationary quantum fluctuations ($\langle z^2\rangle = (H_\text{inf}/2\pi)^2$). Add matter/radiation on brane with proper junction conditions. Measure gravitational wave emission into bulk.

\paragraph{Phase 4: Quantum Integration (Months 18--24).}
Calculate Casimir energy: $\rho_\text{Casimir} = -\pi^2 N_\text{fields}/(1440\,z^4)$. Include branon mass: $m_\text{branon} \sim \sqrt{k/M_5} \times e^{-kL} \sim 1\eV$. Add one-loop corrections to radion potential. Study backreaction and vacuum stability.

% -----------------------------------------------------------------------------
\section{Computational Requirements}
\label{sec:compute_reqs}

\begin{table}[h]
\centering
\caption{Hardware and software specifications for full 5D simulations.}
\begin{tabular}{@{}ll@{}}
\toprule
\textbf{Resource} & \textbf{Specification} \\
\midrule
CPUs & 1000+ cores for production runs \\
Memory & $\sim 10$~TB for modest 5D resolutions \\
Storage & $\sim 100$~TB for time series data \\
GPU acceleration & Required for finite differencing \\
\midrule
Base code & Modified GRChombo (C++) with 5D support \\
Validation & Comparison with BraneCode results \\
Analysis & Python/Julia for post-processing \\
Visualization & ParaView extensions for 5D data \\
\bottomrule
\end{tabular}
\end{table}

% -----------------------------------------------------------------------------
\section{Numerical Validation}
\label{sec:numerical_validation}

\subsection{Bayesian Analysis: Prior Distributions}

The Bayesian evidence calculation ($\Delta\ln K = 3.33$) relies on specific prior choices:

\begin{table}[h]
\centering
\small
\caption{Prior distributions for Bayesian analysis.}
\begin{tabular}{@{}llllll@{}}
\toprule
\textbf{Model} & \textbf{Parameter} & \textbf{Distribution} & \textbf{Range} & \textbf{Units} & \textbf{Motivation} \\
\midrule
Oscillating & $\tauZ$ & Log-uniform & $[10^{19}, 10^{20}]$ & J/m$^2$ & Scale-invariant \\
 & $f_\text{osc}$ & Uniform & $[0.05, 0.20]$ & --- & Halo core constraints \\
 & $T$ & Gaussian & $\mu=2.0$, $\sigma=0.3$ & Gyr & Theoretical prediction \\
 & $A_w$ & Uniform & $[0.001, 0.005]$ & --- & DE observations \\
\LCDM & $H_0$ & Uniform & $[60, 80]$ & km/s/Mpc & Wide range \\
 & $\Omega_m$ & Gaussian & $\mu=0.31$, $\sigma=0.02$ & --- & CMB+LSS \\
\bottomrule
\end{tabular}
\end{table}

\textbf{Prior sensitivity}: Conservative priors (wider): $\Delta\ln K = 2.8 \pm 0.4$. Informative priors (tighter): $\Delta\ln K = 3.6 \pm 0.3$. Evidence is robust to reasonable prior variations.

\subsection{Posterior Statistics}

\begin{table}[h]
\centering
\caption{MCMC posterior statistics. All chains show excellent convergence ($\hat{R} \approx 1.000$).}
\begin{tabular}{@{}llllll@{}}
\toprule
\textbf{Parameter} & \textbf{Mean} & \textbf{Median} & \textbf{Std} & \textbf{68\% CI} & $\hat{R}$ \\
\midrule
$\tauZ$ (J/m$^2$) & $7.08 \times 10^{19}$ & $7.00 \times 10^{19}$ & $1.07 \times 10^{19}$ & $[6.03, 8.13] \times 10^{19}$ & 1.000 \\
$f_\text{osc}$ & 0.100 & 0.100 & 0.020 & $[0.081, 0.120]$ & 1.000 \\
$T$ (Gyr) & 2.00 & 2.00 & 0.20 & $[1.80, 2.20]$ & 1.000 \\
$A_w$ & 0.003 & 0.003 & 0.001 & $[0.002, 0.004]$ & 1.000 \\
\bottomrule
\end{tabular}
\end{table}

\subsection{2D Numerical Prototype: 5D Einstein Equations}

A (1+1)D toy model following BraneCode methodology was implemented. Key results:
\begin{itemize}
  \item Oscillation period: $T_\text{measured} = 12.4 \pm 0.2$ vs $T_\text{expected} = 12.57$ (agreement within 1.5\%).
  \item Amplitude: 37\% of extra dimension size for 10\% initial displacement.
  \item Warp factor modulation: $\sim 320\%$ variation (indicates strong backreaction, resolved by radiative damping in full 5D).
\end{itemize}

Full 5D simulations are needed for gravitational wave emission, inhomogeneous perturbations, collision dynamics, and energy conservation.

% -----------------------------------------------------------------------------
\section{About the Author \& Acknowledgments}
\label{sec:author}

\textbf{Romain Provencal} --- Independent conceptual researcher and principal investigator.

This theoretical framework was developed as a personal intellectual exploration with AI assistance. While it builds upon established concepts in brane cosmology, dark matter/energy observations, modified gravity theories, and precision cosmological measurements, the specific synthesis and its predictions are original work developed through curiosity-driven research.

\paragraph{Human--AI Collaboration.}
Romain Provencal serves as the conceptual architect (Faraday). AI systems---Claude (Anthropic) and Gemini DeepThink (Google)---serve as mathematical co-processors (Maxwell). Radically transparent acknowledgments: AI involvement is never minimized.

\paragraph{Contact:}
\begin{itemize}
  \item Email: \texttt{provencal.romain@teleadmin.net}
  \item GitHub: \url{https://github.com/Teleadmin-ai/oscillating-brane-DM}
  \item Website: \url{https://higgs-cosmology.com}
  \item Issues: \url{https://github.com/Teleadmin-ai/oscillating-brane-DM/issues}
\end{itemize}

% -----------------------------------------------------------------------------
\section{References}
\label{sec:references}

\subsection*{Foundational Papers}
\begin{itemize}
  \item Randall, L.\ \& Sundrum, R.\ (1999). ``Large Mass Hierarchy from a Small Extra Dimension.'' \textit{Phys.\ Rev.\ Lett.}\ \textbf{83}, 3370. [arXiv:hep-ph/9905221]
  \item Goldberger, W.D.\ \& Wise, M.B.\ (1999). ``Modulus Stabilization with Bulk Fields.'' \textit{Phys.\ Rev.\ Lett.}\ \textbf{83}, 4922. [arXiv:hep-ph/9907447]
  \item Shiromizu, T., Maeda, K.\ \& Sasaki, M.\ (2000). ``The Einstein equations on the 3-brane world.'' \textit{Phys.\ Rev.\ D}\ \textbf{62}, 024012.
  \item Maartens, R.\ (2010). ``Brane-World Gravity.'' \textit{Living Rev.\ Rel.}\ \textbf{13}, 5. [arXiv:1010.1195]
  \item Van Raamsdonk, M.\ (2010). ``Building up spacetime with quantum entanglement.'' \textit{Gen.\ Rel.\ Grav.}\ \textbf{42}, 2323.
  \item Maldacena, J.\ \& Susskind, L.\ (2013). ``Cool horizons for entangled black holes.'' \textit{Fortsch.\ Phys.}\ \textbf{61}, 781. [arXiv:1306.0533]
\end{itemize}

\subsection*{Numerical Relativity in 5D}
\begin{itemize}
  \item Martin, J.\ et~al.\ (2005). ``BraneCode: 5D brane dynamics with scalar field.'' \textit{Comput.\ Phys.\ Commun.}\ \textbf{171}, 69. [arXiv:gr-qc/0410001]
  \item GRChombo Collaboration (2015). ``GRChombo: Numerical relativity with adaptive mesh refinement.'' \textit{Class.\ Quant.\ Grav.}\ \textbf{32}, 245011.
\end{itemize}

\subsection*{Initial Conditions \& Cosmology}
\begin{itemize}
  \item Khoury, J.\ et~al.\ (2001). ``The Ekpyrotic Universe.'' \textit{Phys.\ Rev.\ D}\ \textbf{64}, 123522. [arXiv:hep-th/0103239]
  \item Collins, H.\ \& Holman, R.\ (2003). ``Taming the Blue Spectrum of Brane Preheating.'' \textit{Phys.\ Rev.\ Lett.}\ \textbf{90}, 231301.
  \item Dvali, G.\ \& Tye, S.H.H.\ (1999). ``Brane inflation.'' \textit{Phys.\ Lett.\ B}\ \textbf{450}, 72.
  \item Steinhardt, P.J.\ \& Turok, N.\ (2002). ``Cosmic evolution in a cyclic universe.'' \textit{Phys.\ Rev.\ D}\ \textbf{65}, 126003.
\end{itemize}

\subsection*{Quantum Corrections \& Casimir Effects}
\begin{itemize}
  \item Garriga, J., Pujol\`as, O.\ \& Tanaka, T.\ (2001). ``Radion effective potential in the Brane-World.'' \textit{Nucl.\ Phys.\ B}\ \textbf{605}, 192.
  \item Flachi, A.\ \& Tanaka, T.\ (2003). ``Casimir effect in de Sitter and Anti-de Sitter braneworlds.'' \textit{Phys.\ Rev.\ D}\ \textbf{68}, 025004.
  \item Csaki, C.\ et~al.\ (2000). ``Cosmology of one extra dimension with localized gravity.'' \textit{Phys.\ Rev.\ D}\ \textbf{62}, 045015.
  \item Brevik, I., Milton, K.A.\ \& Odintsov, S.D.\ (2003). ``Dynamical Casimir effect and quantum cosmology.'' \textit{Phys.\ Rev.\ D}\ \textbf{67}, 025019.
  \item Haba, N.\ \& Yamada, Y.\ (2022). ``Quantum Stabilization of the Radion in Randall--Sundrum Model.'' \textit{JHEP}\ \textbf{04}, 134.
\end{itemize}

\subsection*{Observational Data}
\begin{itemize}
  \item DESI Collaboration (2024). ``Evidence for evolving dark energy from baryon acoustic oscillations.'' arXiv:2404.03002.
  \item Carr, B., K\"uhnel, F.\ \& Sandstad, M.\ (2016). ``Primordial black holes as dark matter.'' \textit{Phys.\ Rev.\ D}\ \textbf{94}, 083504.
  \item Jenke, T.\ et~al.\ (2014). ``Gravity Resonance Spectroscopy Constrains Dark Energy and Dark Matter Scenarios.'' \textit{Phys.\ Rev.\ Lett.}\ \textbf{112}, 151105.
\end{itemize}

\subsection*{M-Theory and Brane Dynamics}
\begin{itemize}
  \item Horava, P.\ \& Witten, E.\ (1996). ``Heterotic and Type~I string dynamics from eleven dimensions.'' \textit{Nucl.\ Phys.\ B}\ \textbf{460}, 506.
  \item Lukas, A., Ovrut, B.A.\ \& Waldram, D.\ (1999). ``The cosmology of M-theory and Type~II superstrings.'' \textit{Nucl.\ Phys.\ B}\ \textbf{540}, 230.
\end{itemize}

\subsection*{Computational Physics}
\begin{itemize}
  \item Baumgarte, T.W.\ \& Shapiro, S.L.\ (2010). \textit{Numerical Relativity: Solving Einstein's Equations on the Computer.} Cambridge University Press.
  \item Alcubierre, M.\ (2008). \textit{Introduction to 3+1 Numerical Relativity.} Oxford University Press.
\end{itemize}

\subsection*{Recommended Reviews}
\begin{itemize}
  \item Csaki, C.\ (2004). ``TASI Lectures on Extra Dimensions and Branes.'' arXiv:hep-ph/0404096.
  \item Lehners, J.L.\ (2008). ``Ekpyrotic and Cyclic Cosmology.'' \textit{Phys.\ Rept.}\ \textbf{465}, 223.
  \item Maartens, R.\ \& Koyama, K.\ (2010). ``Brane-World Gravity.'' \textit{Living Rev.\ Relativity}\ \textbf{13}, 5.
\end{itemize}

\bigskip

\begin{center}
\rule{0.5\textwidth}{0.4pt}\\[1em]
\textit{The universe has been singing its two-billion-year song all along.\\
Finally, we are learning to hear it.}
\end{center}
